\chapter{Discussion}
\label{chap:discussion}

In the previous chapters we report each experiment on its own terms; in this chapter we read them together: what they establish about the claims of the original work, how the model behaves on synthetic and on real data, how much of each comparison depends on the protocol that selects the hyperparameters, and where the evidence stops.

\section{The Claims of the Original Work}
\label{sec:disc_claims}

The original work reports that \MicroTEDAClus performs on par with established stream clustering algorithms and adapts to concept drift better than them. Under the reassessed methodology, neither claim holds. Against the established benchmark models of the original work, \DenStream, \CluStream and \StreamKM, \MicroTEDAClus is worse than \CluStream and \StreamKM in the Overall ARI of the three static datasets, and at the penultimate batch it is better than none of the three in any dataset and worse than \StreamKM in all six (\cref{sec:reassessment_results}). In two of the three streams with concept drift, its median Steady-State ARI is zero. It learns more slowly than every benchmark model but \AutoCloud, and in the scalability streams of high dimension its microclusters collapse along the stream, which we trace to the dimensionality penalty of the formulation (\cref{subsubsec:dimensionality_penalty}).

The published results were obtained under the conditions that we found in the original methodology (\cref{chap:methodology_reassessment}). The datasets were cleaned of noise, scaled with statistics of the whole stream and ordered by label, which made the clustering of most batches trivial (\cref{sec:datasets}); the evaluation learned each batch before predicting it (\cref{sec:prequential_sequential_offline_evaluation}); the benchmark models ran through layers of interface that inflated their times (\cref{fig:benchmark_java_bridge}); and the implementation that produced the results departs from the formulation of the paper (\cref{sec:model_implementation}). The published comparisons therefore measured these conditions together with the models. The reassessed methodology removes them for every model alike, and under it the original model does not confirm the results published for it.

\section{Synthetic and Real Data}
\label{sec:disc_data}

On the synthetic datasets, \MicroTEDAClusPlus corrected the weaknesses of the original model. It is better than \MicroTEDAClus in the Overall ARI of five of the six datasets and in the Steady-State ARI of all six, faster in both stages and smaller in memory; no benchmark model is better than it in the streams with concept drift, and it is the fastest model of the main experiment (\cref{subsec:mi_main_summary}). Its weakness is the static shapes, where \CluStream and \StreamKM remain better. The modernized experiment, with other benchmark models, longer streams and the selection by rules, moves its position. It is not worse than any benchmark model in the Overall ARI of \BenchmarkOne and \Cassini nor in any scalability stream, but most of them are better than it in \BenchmarkTwo, at the penultimate batch of \RandomRBFGeneratorEvents and in \ChameleonDSThree, and it predicts faster than all of them in every dataset without learning faster (\cref{sec:mod_summary}). On the real datasets, with the same rules for the search ranges, checked again by a pilot, and without its noise label, it is better than every benchmark model in \Shuttle, worse than most of them in four datasets, and left out of \Covertype (\cref{tab:real_quality,tab:real_failures}). These datasets differ from the synthetic ones in what \cref{chap:real_data} names: their drifts come from the conditions of the recording, their classes are not drawn from chosen distributions, and their dimension reaches the hundreds.

\Cref{chap:real_data} also shows where it falls short. Along \InsectsAbrupt, it falls below every benchmark model for a stretch after the first change point, and rises above all of them for a stretch before the third (\cref{fig:real_insects_abrupt}). Its summary grows with the stream: it exceeds the limit in most runs of \Covertype, the longest stream, and in some runs of \Pendigits, whose limit is the smallest, and its peak of memory is the largest of all models in three datasets (\cref{tab:real_failures,tab:real_cost}). In the prediction stage it is never slower than a benchmark model, as on the synthetic data (\cref{tab:real_cost}).

The noise label is a choice that depends on the data. It improved the model on the synthetic datasets, whose streams carry noise (\cref{sec:mi_improvements}), and the real datasets have no noise class, so we run the model without it there (\cref{sec:real_protocol}). The model that no benchmark model surpassed in the synthetic streams with concept drift falls behind several of them in the real ones, the INSECTS datasets, where two or three of the five benchmark models are better (\cref{tab:real_quality}).

\section{The Weight of the Protocol}
\label{sec:disc_protocol}

Every result that we report measures a model as its protocol tunes it, and the protocol weighs on the results as much as some changes to the model do. The probe of \cref{sec:mi_selection} found configurations in the search space of \MicroTEDAClusPlus far better than those that the search selected. The modernized protocol raised the results of the same model in several datasets without any change to it, most of all at the penultimate batch of \BenchmarkOne and \Cassini (\cref{tab:mod_protocols}), and in \ChameleonDSThree it selected finer configurations, of the kind that the probe had found to score better (\cref{subsec:mod_cost}). Its ranges reach past the scales at which a structure summarizes the data, and before each experiment a pilot checks that the selected values do not gather at their ends (\cref{subsec:mod_range_check}). A comparison between models is thus a comparison between models and their selection, which is why we tune every model of a comparison by the same protocol.

The search is also a source of variation. Two searches of a single algorithm give medians that differ without significance (\cref{tab:mi_refinements_tuning_noise}). Differences of that size are not effects of a model, and the tests over many runs are what separate the effects from this variation.

\section{Limits of the Evidence}
\label{sec:disc_limits}

Six limits qualify these results.

\begin{itemize}
    \item \textbf{Streams that repeat.} The Chameleon datasets are fixed files in the main experiment (\cref{subsec:mi_runs}), and the INSECTS datasets and \Covertype keep the order of their files (\cref{subsec:real_order}). Their runs repeat the search and not the stream, so on them the dispersion over the runs is that of the search, and on the Chameleon datasets the validation of \cref{chap:ModelImprovement} confirms the outcome of the search more than that of the model.
    \item \textbf{Number of runs.} The time and the memory of the modernized and the real-data experiments rest on 32 runs (\cref{sec:mod_measurements}). With fewer runs the tests separate fewer models.
    \item \textbf{Models and runs left out.} The limit on the summary leaves \MicroTEDAClusPlus out of \Covertype, which therefore has no comparison, and \BIRCH out of three synthetic datasets, and the failures of \BIRCH leave several other datasets with fewer runs to compare (\cref{tab:mod_failures,tab:real_failures}).
    \item \textbf{Search ranges.} Some ends of the ranges still hold the selected value in a quarter of the runs of a dataset or more, each in one or two datasets: on the synthetic datasets, the upper ends of the number of clusters of \BICO and of the noise distance of \MicroTEDAClusPlus and the lower end of the reassignment ratio of \MiniBatchKMeans, and on the real ones, the lower ends of that ratio and of the cleanup interval of \DBSTREAM (\cref{tab:mod_range_ends,tab:real_range_ends}).
    \item \textbf{Cost.} The execution time and the memory measure implementations and not only algorithms, as \cref{subsec:mi_main_cost} notes and as the reimplementation of \cref{sec:c_implementation} shows for a single algorithm, and the models of a comparison are implemented in different libraries.
    \item \textbf{Datasets.} The synthetic datasets are streams and shapes of two dimensions and scalability streams of up to sixty-four, and the real datasets are eight, chosen by the criteria of \cref{sec:real_datasets}, one of them criticized for its redundant records. Our conclusions about real data extend to data of that kind and not beyond.
\end{itemize}

\clearpage
